\section{数据审计与统一框架}

\subsection{附件结构与质量检查}

程序以只读方式载入附件1至附件4和结果模板，检查工作表、字段、有效行数、时间单调性、重复时刻、缺失值与非数值项。附件1两路分别为3000和4001个观测，采样间隔严格为0.25 s和0.20 s；问题二、三也保留两设备原始事件时刻，不先插值成伪观测。附件4识别出18个射击目标和18个拍照目标。原始输入的 SHA-256 记录在模型合同中，所有正式结果写入项目输出目录。

\subsection{时间偏差符号约定}

时间偏差必须和映射式共同阅读。问题一采用加性修正
\begin{equation}
  t_{2,\mathrm{aligned}}=t_2+\Delta t,
  \label{eq:common-q1-time}
\end{equation}
而问题二、三以方式1为参考，采用
\begin{equation}
  t_{2,\mathrm{ref}}=t_2-\Delta t.
  \label{eq:common-q23-time}
\end{equation}
因此不同问题中同号数值不能脱离公式直接比较。后文同时给出映射方向和物理解释，以免把“设备时钟领先/落后”与“时间轴左移/右移”混淆。

\subsection{稳健配准与异步融合}

多源融合的任务分层与术语采用经典多传感器融合框架\cite{hall1997}，但时间偏差、相对偏差和状态参数全部由本题附件实算。含噪数据采用 Huber 损失降低粗大残差影响\cite{huber1964}。融合阶段不把10 Hz插值点当作观测，而按校正后的4 Hz、5 Hz事件顺序更新常加速度状态；随后使用 RTS 固定区间平滑得到离线轨迹\cite{kalman1960,rts1965}。过程噪声以门控关闭条件下的预检验 NIS 和 Cholesky 白化创新一阶相关选择，避免门控与调参形成反馈。

问题三的偏差判定面对时序相关残差，采用 Newey--West HAC 协方差\cite{neweywest1987}，并用移动区块自助法复核置信区间\cite{kunsch1989}。问题四则将连续可行窗口压缩为候选词典，再以混合整数规划求解；求解器最优状态只证明离散模型，最终仍需回到连续轨迹逐项复核。

\subsection{统一验证原则}

四问均执行三类检查：一是合成数据回收已知参数；二是对搜索范围、插值、统计带宽、过程噪声和扰动尺度作敏感性分析；三是核对正文、参数文件、CSV、图件和工作簿的一致性。凡没有外部真值的量，均表述为“两设备相对差”或“模型内不确定性”，不解释为单一设备的绝对误差。

图\ref{fig:technical-route}给出从原始观测到任务日程的总体技术路线；图\ref{fig:subproblem-flow}进一步标明前问输出和后问输入之间的状态、时钟与不确定性接口。

\begin{figure}[htbp]
  \centering
  \includegraphics[width=.94\textwidth]{paper/technical_route.pdf}
  \caption{配准、融合、偏差判定与任务调度的总体技术路线}
  \label{fig:technical-route}
\end{figure}
\begin{figure}[htbp]
  \centering
  \includegraphics[width=.88\textwidth]{paper/subproblem_flow.pdf}
  \caption{四个子问题的依赖关系与共享数据接口}
  \label{fig:subproblem-flow}
\end{figure}
